\chapterbackmatter{Solutions to Supplementary Exercises}

\begin{enumerate}
\SupplementarySolution{1}{The Proca Field}
\begin{enumerate}
\item
The Maxwell term is invariant under
\(A_\mu\mapsto A_\mu+\partial_\mu\Lambda\), but the mass term changes by
\[
 -\frac12m^2\left[
 2A^\mu\partial_\mu\Lambda
 +\partial^\mu\Lambda\,\partial_\mu\Lambda\right],
\]
which is not a total derivative for a general local \(\Lambda\).  Varying
\(A_\nu\) gives
\[
 \partial_\mu F^{\mu\nu}-m^2A^\nu=0.
\]
Taking a divergence and using the antisymmetry of \(F^{\mu\nu}\) yields
\[
 m^2\partial_\nu A^\nu=0.
\]
Since \(m^2>0\), the field is transverse.  Expanding the remaining
equation then gives
\[
 \boxed{\partial_\mu A^\mu=0,\qquad
 (\Box-m^2)A^\nu=0.}
\]

\item
A metric variation gives the symmetric energy-momentum tensor
\[
 T_{\mu\nu}
 =F_{\mu\rho}F_\nu{}^\rho
 -\frac14\eta_{\mu\nu}F_{\rho\sigma}F^{\rho\sigma}
 +m^2\left(A_\mu A_\nu-\frac12\eta_{\mu\nu}A_\rho A^\rho\right).
\]
The canonical Noether tensor differs from this expression, on the equations
of motion, by the divergence of a local superpotential.  Thus their
integrated energies differ only by a spatial boundary term.  With
\(E_i=F_{0i}\) and \(B_i=\frac12\epsilon_{ijk}F_{jk}\),
\[
 \boxed{
 T_{00}
 =\frac12\left(\mathbf E^2+\mathbf B^2\right)
 +\frac{m^2}{2}\left[(A^0)^2+\mathbf A^2\right]\geq0.}
\]
This is manifestly positive; the qualification about a total spatial
derivative is needed only if one starts from the canonical tensor.

\item
For
\[
 A^\mu(x)=\epsilon^\mu(p)e^{-ip\cdot x},
\]
the field equations become
\[
 p^2=-m^2,\qquad p_\mu\epsilon^\mu(p)=0.
\]
In the rest frame \(p^\mu=(m,\mathbf0)\), transversality sets
\(\epsilon^0=0\), leaving three independent spatial vectors.  For general
\(\mathbf p\), choose two orthonormal vectors
\(\mathbf e_1,\mathbf e_2\perp\mathbf p\) and take
\[
 \epsilon^\mu_{1,2}=(0,\mathbf e_{1,2}),\qquad
 \epsilon^\mu_L
 =\frac1m\left(|\mathbf p|,p^0\widehat{\mathbf p}\right).
\]
All three have unit spacelike norm and are orthogonal to \(p^\mu\).  Their
completeness relation is
\[
 \sum_{\lambda=1}^{3}
 \epsilon_\mu(p,\lambda)\epsilon_\nu(p,\lambda)^*
 =\eta_{\mu\nu}+\frac{p_\mu p_\nu}{m^2}.
\]

\item
Before imposing the constraint, the differential operator is
\[
 \Pi_{\mu\nu}
 =(\Box-m^2)\eta_{\mu\nu}-\partial_\mu\partial_\nu .
\]
In momentum space it becomes
\[
 \Pi_{\mu\nu}(p)
 =-(p^2+m^2)\eta_{\mu\nu}+p_\mu p_\nu .
\]
Its inverse is
\[
 [\Pi^{-1}(p)]_{\mu\nu}
 =-\frac{\eta_{\mu\nu}+p_\mu p_\nu/m^2}{p^2+m^2},
\]
as direct multiplication verifies.  The Feynman prescription therefore
gives
\[
 \boxed{
 \Delta_{\mu\nu}(x-y)
 =\int\frac{d^4p}{(2\pi)^4}\,
 \frac{-i\left(\eta_{\mu\nu}+p_\mu p_\nu/m^2\right)}
 {p^2+m^2-i0}\,
 e^{ip\cdot(x-y)} .}
\]
The numerator is precisely the polarization sum found in part (c).
\end{enumerate}

\SupplementarySolution{2}{Massless QED with a Pauli Interaction}
For brevity define
\[
 C^{\mu\nu}=[\gamma^\mu,\gamma^\nu],
 \qquad B^{\mu\nu}=\bar\psi C^{\mu\nu}\psi.
\]
\begin{enumerate}
\item
In four dimensions,
\([psi]=3/2\), \([A_\mu]=1\), and \([F_{\mu\nu}]=2\).  Hence the
minimal operator \(\bar\psi\gamma^\mu A_\mu\psi\) has dimension four,
whereas \(B^{\mu\nu}F_{\mu\nu}\) has dimension five.  Therefore
\[
 \boxed{[e]=0\quad\hbox{(marginal)},\qquad
 [\lambda]=-1\quad\hbox{(irrelevant)}.}
\]
This is the classical or engineering classification; quantum running can
make the marginal coupling logarithmically scale dependent.
\item
For \(x'=\Lambda x\),
\[
 \psi'(x')=S(\Lambda)\psi(x),\qquad
 \bar\psi'(x')=\bar\psi(x)S(\Lambda)^{-1},\qquad
 A'^\mu(x')=\Lambda^\mu{}_{\nu}A^\nu(x),
\]
where
\(S^{-1}\gamma^\mu S=\Lambda^\mu{}_{\nu}\gamma^\nu\).
Thus \(D_\mu\psi\) is a vector-spinor, \(F_{\mu\nu}\) is an
antisymmetric tensor, and \(\bar\psi C^{\mu\nu}\psi\) is an
antisymmetric tensor.  Every index in the Lagrangian is contracted.

The volume gamma matrices obey
\[
 \beta^{-1}(\gamma^\mu)^\dagger\beta=-\gamma^\mu,
 \qquad
 \beta^{-1}(C^{\mu\nu})^\dagger\beta=-C^{\mu\nu}.
\]
The Maxwell term is real, \(iB^{\mu\nu}F_{\mu\nu}\) is Hermitian, and
the gauge interaction is Hermitian because
\(\bar\psi\gamma^\mu\psi\) is anti-Hermitian.  Finally,
\[
 \left(-\bar\psi\gamma^\mu\partial_\mu\psi\right)^\dagger
 =-\bar\psi\gamma^\mu\partial_\mu\psi
 +\partial_\mu(\bar\psi\gamma^\mu\psi).
\]
Hence the action is real; the only mismatch in the density is the displayed
total derivative.
\item
Independent variation of \(\bar\psi\) and \(\psi\) gives
\[
 \boxed{
 \gamma^\mu D_\mu\psi-i\lambda C^{\mu\nu}F_{\mu\nu}\psi=0,}
\]
\[
 \boxed{
 (D_\mu\bar\psi)\gamma^\mu
 +i\lambda\bar\psi C^{\mu\nu}F_{\mu\nu}=0,}
 \qquad
 D_\mu\bar\psi=\partial_\mu\bar\psi-ieA_\mu\bar\psi.
\]
Variation of the gauge field, including the derivative inside
\(F_{\mu\nu}\), yields
\[
 \boxed{
 \partial_\mu F^{\mu\nu}
 =ie\bar\psi\gamma^\nu\psi
 +2i\lambda\partial_\mu
 \left(\bar\psi C^{\mu\nu}\psi\right).}
\]
The factor two follows from the antisymmetry of \(C^{\mu\nu}\).
\item
For an arbitrary real function \(\Lambda(x)\),
\[
 A'_\mu=A_\mu+\partial_\mu\Lambda,
 \qquad
 \psi'=e^{-ie\Lambda}\psi,
 \qquad
 \bar\psi'=\bar\psi e^{ie\Lambda}.
\]
Then \(D'_\mu\psi'=e^{-ie\Lambda}D_\mu\psi\), while
\(F'_{\mu\nu}=F_{\mu\nu}\); all three terms are locally invariant.
Taking \(\Lambda\) constant gives the global symmetry.  With the
unit-normalized infinitesimal transformation
\(\delta\psi=-i\epsilon\psi\), Noether's theorem gives
\[
 \boxed{j^\mu=i\bar\psi\gamma^\mu\psi,qquad
 Q=\int d^3\mathbf x\,j^0.}
\]
Its divergence is
\[
 \partial_\mu j^\mu
 =i\left[(D_\mu\bar\psi)\gamma^\mu\psi
 +\bar\psi\gamma^\mu D_\mu\psi\right]=0,
\]
because the two Pauli terms supplied by the conjugate field equations
cancel.  Thus \(\dot Q=0\) when the spatial flux vanishes at infinity.

The current sourcing Maxwell's equation is
\[
 J_{\rm Max}^\nu
 =e j^\nu+2i\lambda\partial_\mu B^{\mu\nu}.
\]
The second term is an identically conserved magnetization current:
\(\partial_\nu\partial_\mu B^{\mu\nu}=0\).  Its contribution to the
integrated charge is a spatial boundary term, so for localized fields
\(\int d^3\mathbf x\,J_{\rm Max}^0=eQ\), even though the local source
density is not simply \(e j^0\).
\item
Because \(\{\gamma_5,\gamma^\mu\}=0\), the massless kinetic and minimal
coupling terms are invariant under a constant axial rotation; the adjoint
transforms as \(\bar\psi\mapsto\bar\psi e^{i\alpha\gamma_5}\).  In
contrast, \([\gamma_5,C^{\mu\nu}]=0\), so
\[
 \bar\psi C^{\mu\nu}\psi
 \longmapsto
 \bar\psi C^{\mu\nu}e^{2i\alpha\gamma_5}\psi.
\]
For a continuous chiral symmetry one must therefore set
\[
 \boxed{\lambda=0.}
\]
If a mass term is present, \(\bar\psi\psi\) acquires the same
\(e^{2i\alpha\gamma_5}\) factor, so one also needs \(m=0\).

For \(\alpha=\alpha(x)\), the kinetic term produces
\[
 \delta\mathcal L
 =-i(\partial_\mu\alpha)
 \bar\psi\gamma^\mu\gamma_5\psi.
\]
A variation of the existing vector gauge field couples instead to
\(\bar\psi\gamma^\mu\psi\); adjusting \(e\) or \(\lambda\) cannot cancel
the axial current.  A new axial gauge field would be required, so no
choice of the two displayed couplings makes this chiral symmetry local.
\end{enumerate}

\SupplementarySolution{3}{Electron--Positron Annihilation to Muons}
\begin{enumerate}
\item
There is one tree diagram: the electron--positron pair annihilates into a
virtual photon, which creates the muon pair.  Suppressing a
convention-dependent common phase, the amplitude is
\[
 \boxed{
 \mathcal M
 =\frac{e^2}{s}
 \left[\bar v_e(p_{e^+})\gamma^\mu u_e(p_{e^-})\right]
 \left[\bar u_\mu(p_{\mu^-})\gamma_\mu
 v_\mu(p_{\mu^+})\right].}
\]

\item
The standard spin sums, translated to the volume gamma convention, reduce
the product of currents to two Dirac traces.  Writing
\[
 A_t=m_e^2+m_\mu^2-t,\qquad
 A_u=m_e^2+m_\mu^2-u,
\]
the result is
\[
 \boxed{
 \overline{|\mathcal M|^2}
 =\frac{2e^4}{s^2}
 \left[
 A_t^2+A_u^2+2s(m_e^2+m_\mu^2)
 \right].}
\]
Using the four-gamma trace identity and
\(s+t+u=2m_e^2+2m_\mu^2\) gives the boxed expression.  It reduces to
\(2e^4(t^2+u^2)/s^2\) when both masses vanish.

\item
Define
\[
 \beta_e=\sqrt{1-\frac{4m_e^2}{s}},\qquad
 \beta_\mu=\sqrt{1-\frac{4m_\mu^2}{s}}.
\]
In the center-of-mass frame,
\[
 |\mathbf p_i|=\frac{\sqrt s}{2}\beta_e,\qquad
 |\mathbf p_f|=\frac{\sqrt s}{2}\beta_\mu,
\]
and
\[
 \begin{split}
 t&=m_e^2+m_\mu^2-\frac{s}{2}
 +\frac{s}{2}\beta_e\beta_\mu\cos\theta,\\
 u&=m_e^2+m_\mu^2-\frac{s}{2}
 -\frac{s}{2}\beta_e\beta_\mu\cos\theta.
 \end{split}
\]
Substitution gives
\[
 \boxed{
 \overline{|\mathcal M|^2}
 =e^4\left[
 3-\beta_e^2-\beta_\mu^2
 +\beta_e^2\beta_\mu^2\cos^2\theta
 \right].}
\]

\item
The supplied two-body formula now yields
\[
 \boxed{
 \frac{d\sigma}{d\Omega}
 =\frac{\alpha^2}{4s}\frac{\beta_\mu}{\beta_e}
 \left[
 3-\beta_e^2-\beta_\mu^2
 +\beta_e^2\beta_\mu^2\cos^2\theta
 \right],}
\qquad \alpha=\frac{e^2}{4\pi}.
\]
Using \(\int d\Omega\,\cos^2\theta=4\pi/3\), the total cross section is
\[
 \boxed{
 \sigma(e^+e^-\to\mu^+\mu^-)
 =\frac{\pi\alpha^2}{s}\frac{\beta_\mu}{\beta_e}
 \left[
 3-\beta_e^2-\beta_\mu^2
 +\frac13\beta_e^2\beta_\mu^2
 \right].}
\]
As a check, when \(m_e\to0\),
\[
 \sigma
 =\frac{4\pi\alpha^2}{3s}\,
 \beta_\mu\left(1+\frac{2m_\mu^2}{s}\right),
\]
the usual result.  No identical-particle factor is present because
\(\mu^+\) and \(\mu^-\) are distinct.
\end{enumerate}

\SupplementarySolution{4}{Gauge Symmetry, QED Rules, and Three Tree Processes}
\begin{enumerate}
\item
A gauge symmetry is a local redundancy in the field variables: functions
related by the gauge transformation describe the same physical
configuration.  For an electron of charge \(-e\), define
\[
 D_\mu=\partial_\mu+ieA_\mu,\qquad
 \mathcal L_{\rm QED}
 =-\frac14F_{\mu\nu}F^{\mu\nu}
 -\bar\Psi(\gamma^\mu D_\mu+m)\Psi .
\]
Under
\[
 A'_\mu=A_\mu+\partial_\mu\Lambda,\qquad
 \Psi'=e^{-ie\Lambda}\Psi,\qquad
 \bar\Psi'=\bar\Psi e^{ie\Lambda},
\]
one has
\[
 D'_\mu\Psi'=e^{-ie\Lambda}D_\mu\Psi,
 \qquad F'_{\mu\nu}=F_{\mu\nu}.
\]
The phases cancel in the Dirac bilinear, and the Maxwell term is unchanged,
which proves the local invariance.

\pagebreak[3]
\item
Suppressing the momentum-conserving delta functions and powers of \(2\pi\),
a consistent set of rules in the volume convention is as follows.  Direct
each fermion line along its arrow and use
\[
 \begin{split}
 &\text{vertex:}\quad e\gamma^\mu,\\
 &\text{fermion line:}\quad
 \frac{-i[-i\slashed q+m]}{q^2+m^2-i0},\\
 &\text{photon line:}\quad
 \frac{-i}{q^2-i0}
 \left[\eta_{\mu\nu}
+(\xi-1)\frac{q_\mu q_\nu}{q^2-i0}\right].
 \end{split}
\]
Feynman gauge is \(\xi=1\).  Attach \(u,\bar u\) to incoming and outgoing
electrons, \(\bar v,v\) to incoming and outgoing positrons, and
\(\epsilon_\mu,\epsilon_\mu^*\) to incoming and outgoing photons,
respectively.  Integrate every undetermined internal momentum.  Include a
minus sign for each closed fermion loop and for every odd permutation needed
to restore a fixed ordering of external fermionic operators.  These rules,
together with four-momentum conservation at each vertex, are complete for
the requested tree graphs.

\item
Each process has two tree topologies.  The labels shown are also the
complete one-particle quantum numbers required in the question.

% Photons are drawn wavy and charged lines solid, matching the convention used
% for the self-energy graphs in the Chapter 11 problems.  Fermion-flow arrows
% sit at the middle of each segment rather than at its end, so no arrowhead
% lands on a vertex, and every momentum label hangs off the free end of its
% own leg instead of straddling it.
\smallskip
\noindent\emph{Bhabha scattering: \(s\)-channel annihilation and
\(t\)-channel exchange}
\begin{center}
\begin{tikzpicture}[x=0.72cm,y=0.72cm,baseline=-0.5ex,line width=0.65pt,
  ferm/.style={postaction={decorate},decoration={markings,
    mark=at position #1 with {\arrow{Stealth[length=4.6pt,width=3.4pt]}}}},
  ferm/.default=0.55,
  photon/.style={decorate,
    decoration={snake,amplitude=1.05pt,segment length=4.6pt}},
  vertex/.style={circle,fill,inner sep=1.25pt},
  every node/.style={font=\small}]
 \coordinate (a) at (-1.05,0); \coordinate (b) at (1.05,0);
 \draw[ferm] (-2.45,1.15)--(a);
 \draw[ferm] (a)--(-2.45,-1.15);
 \draw[photon] (a)--(b);
 \draw[ferm] (b)--(2.45,1.15);
 \draw[ferm] (2.45,-1.15)--(b);
 \node[vertex] at (a) {}; \node[vertex] at (b) {};
 \node[anchor=south east] at (-2.5,1.2) {\(\scriptstyle p_1,s_1\)};
 \node[anchor=north east] at (-2.5,-1.2) {\(\scriptstyle p_2,s_2\)};
 \node[anchor=south west] at (2.5,1.2) {\(\scriptstyle p_3,s_3\)};
 \node[anchor=north west] at (2.5,-1.2) {\(\scriptstyle p_4,s_4\)};
 \node[above=3pt] at (0,0) {\(\scriptstyle\gamma^*(s)\)};
\end{tikzpicture}
\qquad
\begin{tikzpicture}[x=0.72cm,y=0.72cm,baseline=-0.5ex,line width=0.65pt,
  ferm/.style={postaction={decorate},decoration={markings,
    mark=at position #1 with {\arrow{Stealth[length=4.6pt,width=3.4pt]}}}},
  ferm/.default=0.55,
  photon/.style={decorate,
    decoration={snake,amplitude=1.05pt,segment length=4.6pt}},
  vertex/.style={circle,fill,inner sep=1.25pt},
  every node/.style={font=\small}]
 \coordinate (a) at (0,0.9); \coordinate (b) at (0,-0.9);
 \draw[ferm] (-2.3,0.9)--(a);
 \draw[ferm] (a)--(2.3,0.9);
 \draw[ferm] (b)--(-2.3,-0.9);
 \draw[ferm] (2.3,-0.9)--(b);
 \draw[photon] (a)--(b);
 \node[vertex] at (a) {}; \node[vertex] at (b) {};
 \node[above] at (-1.2,0.95) {\(\scriptstyle p_1,s_1\)};
 \node[above] at (1.2,0.95) {\(\scriptstyle p_3,s_3\)};
 \node[below] at (-1.2,-0.95) {\(\scriptstyle p_2,s_2\)};
 \node[below] at (1.2,-0.95) {\(\scriptstyle p_4,s_4\)};
 \node[right=4pt] at (0,0) {\(\scriptstyle\gamma^*(t)\)};
\end{tikzpicture}
\end{center}

\noindent\emph{Pair annihilation: the two orders of photon emission}
\begin{center}
\begin{tikzpicture}[x=0.72cm,y=0.72cm,baseline=-0.5ex,line width=0.65pt,
  ferm/.style={postaction={decorate},decoration={markings,
    mark=at position #1 with {\arrow{Stealth[length=4.6pt,width=3.4pt]}}}},
  ferm/.default=0.55,
  photon/.style={decorate,
    decoration={snake,amplitude=1.05pt,segment length=4.6pt}},
  vertex/.style={circle,fill,inner sep=1.25pt},
  every node/.style={font=\small}]
 \coordinate (a) at (0,0.78); \coordinate (b) at (0,-0.78);
 \draw[ferm] (-2.0,1.6)--(a);
 \draw[ferm] (a)--(b);
 \draw[ferm] (b)--(-2.0,-1.6);
 \draw[photon] (a)--(2.0,1.6);
 \draw[photon] (b)--(2.0,-1.6);
 \node[vertex] at (a) {}; \node[vertex] at (b) {};
 \node[anchor=south east] at (-2.05,1.65) {\(\scriptstyle p_1,s_1\)};
 \node[anchor=north east] at (-2.05,-1.65) {\(\scriptstyle p_2,s_2\)};
 \node[anchor=south west] at (2.05,1.65) {\(\scriptstyle k_1,\lambda_1\)};
 \node[anchor=north west] at (2.05,-1.65) {\(\scriptstyle k_2,\lambda_2\)};
\end{tikzpicture}
\qquad
\begin{tikzpicture}[x=0.72cm,y=0.72cm,baseline=-0.5ex,line width=0.65pt,
  ferm/.style={postaction={decorate},decoration={markings,
    mark=at position #1 with {\arrow{Stealth[length=4.6pt,width=3.4pt]}}}},
  ferm/.default=0.55,
  photon/.style={decorate,
    decoration={snake,amplitude=1.05pt,segment length=4.6pt}},
  vertex/.style={circle,fill,inner sep=1.25pt},
  every node/.style={font=\small}]
 \coordinate (a) at (0,0.78); \coordinate (b) at (0,-0.78);
 \draw[ferm] (-2.0,1.6)--(a);
 \draw[ferm] (a)--(b);
 \draw[ferm] (b)--(-2.0,-1.6);
 \draw[photon] (a)--(2.0,1.6);
 \draw[photon] (b)--(2.0,-1.6);
 \node[vertex] at (a) {}; \node[vertex] at (b) {};
 \node[anchor=south east] at (-2.05,1.65) {\(\scriptstyle p_1,s_1\)};
 \node[anchor=north east] at (-2.05,-1.65) {\(\scriptstyle p_2,s_2\)};
 \node[anchor=south west] at (2.05,1.65) {\(\scriptstyle k_2,\lambda_2\)};
 \node[anchor=north west] at (2.05,-1.65) {\(\scriptstyle k_1,\lambda_1\)};
\end{tikzpicture}
\end{center}

\noindent\emph{Positron Compton scattering: the two orders along the
positron line}
\begin{center}
\begin{tikzpicture}[x=0.72cm,y=0.72cm,baseline=-0.5ex,line width=0.65pt,
  ferm/.style={postaction={decorate},decoration={markings,
    mark=at position #1 with {\arrow{Stealth[length=4.6pt,width=3.4pt]}}}},
  ferm/.default=0.55,
  photon/.style={decorate,
    decoration={snake,amplitude=1.05pt,segment length=4.6pt}},
  vertex/.style={circle,fill,inner sep=1.25pt},
  every node/.style={font=\small}]
 \coordinate (a) at (-0.95,0); \coordinate (b) at (0.95,0);
 \draw[ferm] (a)--(-2.55,0);
 \draw[ferm] (b)--(a);
 \draw[ferm] (2.55,0)--(b);
 \draw[photon] (-1.95,1.75)--(a);
 \draw[photon] (b)--(1.95,1.75);
 \node[vertex] at (a) {}; \node[vertex] at (b) {};
 \node[anchor=east] at (-2.6,0) {\(\scriptstyle p_1,s_1\)};
 \node[anchor=west] at (2.6,0) {\(\scriptstyle p_2,s_2\)};
 \node[anchor=south] at (-1.95,1.82) {\(\scriptstyle k_1,\lambda_1\)};
 \node[anchor=south] at (1.95,1.82) {\(\scriptstyle k_2,\lambda_2\)};
 \path (0,-0.9);
\end{tikzpicture}
\qquad
\begin{tikzpicture}[x=0.72cm,y=0.72cm,baseline=-0.5ex,line width=0.65pt,
  ferm/.style={postaction={decorate},decoration={markings,
    mark=at position #1 with {\arrow{Stealth[length=4.6pt,width=3.4pt]}}}},
  ferm/.default=0.55,
  photon/.style={decorate,
    decoration={snake,amplitude=1.05pt,segment length=4.6pt}},
  vertex/.style={circle,fill,inner sep=1.25pt},
  every node/.style={font=\small}]
 \coordinate (a) at (-0.95,0); \coordinate (b) at (0.95,0);
 \draw[ferm] (a)--(-2.55,0);
 \draw[ferm] (b)--(a);
 \draw[ferm] (2.55,0)--(b);
 \draw[photon] (-1.95,1.75)--(a);
 \draw[photon] (b)--(1.95,1.75);
 \node[vertex] at (a) {}; \node[vertex] at (b) {};
 \node[anchor=east] at (-2.6,0) {\(\scriptstyle p_1,s_1\)};
 \node[anchor=west] at (2.6,0) {\(\scriptstyle p_2,s_2\)};
 \node[anchor=south] at (-1.95,1.82) {\(\scriptstyle k_2,\lambda_2\)};
 \node[anchor=south] at (1.95,1.82) {\(\scriptstyle k_1,\lambda_1\)};
 \path (0,-0.9);
\end{tikzpicture}
\end{center}

\item
It is convenient to strip one common overall phase and define
\[
 G(q)\equiv\frac{-i\slashed q+m}{q^2+m^2-i0},
 \qquad
 \slashed\epsilon\equiv\gamma^\mu\epsilon_\mu .
\]
For Bhabha scattering, a fixed ordering of the four external fermion
operators gives
\[
 \begin{split}
 \mathcal A_{\rm Bh}=e^2\Bigg[
 &\frac{
 [\bar v_2\gamma^\mu u_1]
 [\bar u_3\gamma_\mu v_4]}
 {(p_1+p_2)^2-i0}\\
 &-\frac{
 [\bar u_3\gamma^\mu u_1]
 [\bar v_2\gamma_\mu v_4]}
 {(p_1-p_3)^2-i0}
 \Bigg].
 \end{split}
\]
The relative minus sign is the fermion-interchange sign; a different common
ordering changes only an overall sign.

For pair annihilation, \(p_1+p_2=k_1+k_2\), and the two possible photon
orders give
\[
 \boxed{
 \mathcal A_{\gamma\gamma}
 =e^2\bar v_2\left[
 \slashed\epsilon_2^*G(p_1-k_1)\slashed\epsilon_1^*
 +\slashed\epsilon_1^*G(p_1-k_2)\slashed\epsilon_2^*
 \right]u_1 .}
\]
The expression is symmetric under interchange of the two final photons, as
it must be.

Finally \(p_1+k_1=p_2+k_2\) for positron Compton scattering.  Following
momentum along the fermion arrow gives
\[
 \boxed{
 \mathcal A_{e^+\gamma}
 =e^2\bar v_1\left[
 \slashed\epsilon_1G(-p_1-k_1)\slashed\epsilon_2^*
 +\slashed\epsilon_2^*G(-p_1+k_2)\slashed\epsilon_1
 \right]v_2 .}
\]
In all three answers the omitted common phase depends only on the convention
used to define \(S=\mathbf1+iT\).  Replacing either external photon
polarization by its momentum makes the two terms cancel after the external
Dirac equations are used, providing a check on the relative signs and
momentum routing.
\end{enumerate}

\SupplementarySolution{5}{Canonical Structure of Scalar Electrodynamics}
Write
\[
 D_\mu\Phi=(\partial_\mu-iqA_\mu)\Phi,\qquad
 (D_\mu\Phi)^*=(\partial_\mu+iqA_\mu)\Phi^* .
\]
\begin{enumerate}
\item
Variation with respect to \(A_\mu\) gives
\[
 \partial_\nu F^{\nu\mu}=j^\mu,\qquad
 \boxed{
 j^\mu=iq\left[
 \Phi^*D^\mu\Phi-(D^\mu\Phi)^*\Phi
 \right].}
\]
Equivalently,
\[
 j^\mu
 =iq\left(\Phi^*\partial^\mu\Phi
 -\partial^\mu\Phi^*\,\Phi\right)
 +2q^2A^\mu|\Phi|^2 .
\]
Thus the term quadratic in \(q\), which would be absent for a fixed
background current, is required by the dynamical charged field.

\item
Varying \(\Phi^*\) and integrating the covariant derivative by parts gives
\[
 \boxed{(D_\mu D^\mu-m^2)\Phi=0.}
\]
The conjugate equation is
\((D_\mu^*D^{*\mu}-m^2)\Phi^*=0\).  Multiplying these equations by
\(\Phi^*\) and \(\Phi\), subtracting, and using the neutrality of \(j^\mu\)
also verifies \(\partial_\mu j^\mu=0\).

\item
Because the mostly-plus scalar kinetic term is
\((D_0\Phi)^*D_0\Phi-(\mathbf D\Phi)^*\cdot\mathbf D\Phi\), the canonical
momenta are
\[
 \boxed{\begin{aligned}
 \Pi_\Phi
 &=\frac{\partial\mathcal L}{\partial\dot\Phi}
 =(D_0\Phi)^*=\dot\Phi^*+iqA_0\Phi^*,\\
 \Pi_{\Phi^*}
 &=\frac{\partial\mathcal L}{\partial\dot\Phi^*}
 =D_0\Phi=\dot\Phi-iqA_0\Phi .
 \end{aligned}}
\]
The charge density and spatial current are therefore
\[
 \boxed{
 j^0=iq\left(\Pi_\Phi\Phi-\Phi^*\Pi_{\Phi^*}\right),}
\]
\[
 \boxed{
 \mathbf j
 =iq\left[\Phi^*\boldsymbol\nabla\Phi
 -(\boldsymbol\nabla\Phi^*)\Phi\right]
 +2q^2\mathbf A|\Phi|^2 .}
\]
These expressions contain only the canonical coordinates, their conjugate
momenta, \(A_i\), and spatial derivatives; \(A_0\) has disappeared from the
charge density.

\item
Under the stated transformation,
\[
 D'_\mu\Phi'
 =(\partial_\mu-iqA_\mu-iq\partial_\mu\Omega)
   e^{iq\Omega}\Phi
 =e^{iq\Omega}D_\mu\Phi .
\]
Its conjugate transforms with the inverse phase, while
\[
 F'_{\mu\nu}
 =F_{\mu\nu}
 +\partial_\mu\partial_\nu\Omega
 -\partial_\nu\partial_\mu\Omega
 =F_{\mu\nu}.
\]
Hence each of the Maxwell, scalar kinetic, and mass terms is separately
invariant.  The canonical variables transform as
\[
 \Phi'\!=e^{iq\Omega}\Phi,\quad
 \Pi'_\Phi=e^{-iq\Omega}\Pi_\Phi,\quad
 \Phi'^*\!=e^{-iq\Omega}\Phi^*,\quad
 \Pi'_{\Phi^*}=e^{iq\Omega}\Pi_{\Phi^*},
\]
so the phase-space form of \(j^\mu\) is gauge invariant as well.
\end{enumerate}

\SupplementarySolution{6}{Magnetic Flux, Holonomy, and Dirac Quantization}
The field strength is the intrinsic two-form
\[
 F=dA=(\partial_\theta A_\varphi)\,d\theta\wedge d\varphi
 =\frac{n}{2e}\sin\theta\,d\theta\wedge d\varphi.
\]
It is proportional to the area form on the unit sphere, so the radial
magnetic flux density is constant.  Its total flux is
\[
 \boxed{
 \int_{S^2}F
 =\frac{n}{2e}\int_0^\pi d\theta\,\sin\theta
 \int_0^{2\pi}d\varphi
 =\frac{2\pi n}{e}.}
\]
Along a latitude, \(A_\varphi\) is constant, and hence
\[
 \boxed{h(\theta)=-\frac{\pi n}{e}(1+\cos\theta).}
\]
At the south pole, \(\theta=\pi\), this is zero.  At the north pole,
\(\theta=0\), the coordinate circle also shrinks to a point but
\[
 h(0)=-\frac{2\pi n}{e}.
\]
The gauge potential is singular in this patch at that pole; the
gauge-invariant Wilson line must nevertheless become the identity:
\[
 W(0)=e^{ieh(0)}=e^{-2\pi i n}=1.
\]
Therefore
\[
 \boxed{n\in\mathbb Z.}
\]
Equivalently, the magnetic flux is quantized in units of \(2\pi/e\).
The apparent nonzero holonomy of a vanishing contour records the
nontrivial transition function between gauge patches, not a physical
singularity of the field strength.

\SupplementarySolution{7}{The Gupta--Bleuler Physical Space and Residual Gauge Modes}
\begin{enumerate}
\item
In Feynman gauge the four components behave as oscillators with an
indefinite metric.  The canonical equal-time algebra implies
\[
 [A_0(t,\mathbf x),\partial_\mu A^\mu(t,\mathbf y)]
 =\pm i\delta^3(\mathbf x-\mathbf y),
\]
where the displayed sign depends on whether the time index is raised before
defining the canonical momentum.  The crucial point is that the right-hand
side is nonzero.  If \(\partial_\mu A^\mu\) were the zero operator, its
commutator with every operator would vanish, contradicting this canonical
relation.  Lorenz gauge must instead be imposed weakly as a condition on
physical states.
\item
Decompose the divergence into annihilation and creation pieces,
\[
 \partial_\mu A^\mu
 =(\partial\cdot A)_{\rm ann}+(\partial\cdot A)_{\rm cre}.
\]
The Gupta--Bleuler condition is
\[
 (\partial\cdot A)_{\rm ann}\ket{\psi_{\rm phys}}=0.
\]
For each momentum, the divergence contains the combination
\[
 b_+(\mathbf k)=\frac{a_L(\mathbf k)+a_0(\mathbf k)}{\sqrt2}
\]
up to a harmless convention-dependent sign, so the condition is
\(b_+(\mathbf k)\ket{\psi_{\rm phys}}=0\).  It eliminates negative-norm
excitations without requiring the creation part to annihilate the state.
Because every annihilation operator kills the vacuum, the vacuum remains
physical.  Matrix elements of the full divergence between physical states
then vanish.
\item
Normalize the longitudinal and time-like modes by
\[
 [a_L,a_L^\dagger]=1,\qquad [a_0,a_0^\dagger]=-1,
\]
and define \(b_\pm=(a_L\pm a_0)/\sqrt2\).  Direct expansion gives
\[
 [b_+,b_+^\dagger]=[b_-,b_-^\dagger]=0,\qquad
 [b_+,b_-^\dagger]=[b_-,b_+^\dagger]=1,
\]
with momentum delta functions understood.  Since the physical condition is
\(b_+\ket{\psi}=0\), a \(b_-^\dagger\) excitation is forbidden:
commuting \(b_+\) through it leaves the state below.  In contrast,
\(b_+^\dagger\) excitations are allowed because their self-commutator is
zero.  They also have zero norm and are the null photons.
\item
The oscillator algebra shows that a state satisfying
\(b_+\ket{\psi}=0\) can be built from transverse creation operators and
\(b_+^\dagger\), but not from \(b_-^\dagger\).  Move all
\(b_+^\dagger\) factors to one side and write
\[
 \ket{\psi}=\ket{\psi_T}+\ket{\chi},
\]
where every term in \(\ket{\chi}\) contains at least one
\(b_+^\dagger\).  Its inner product with any physical state vanishes:
after Hermitian conjugation, a \(b_+\) can be commuted through all
transverse and \(b_+^\dagger\) operators without producing a number, and
then annihilates the vacuum.  Thus the null vectors form a radical of the
inner product.  The quotient
\[
 \mathcal H_{\rm phys}=\mathcal H_{\rm small}/\mathcal N
\]
is represented uniquely by transverse states and has a positive-definite
inner product with exactly two photon polarizations.
\item
For a one-null-photon state, only the annihilation part of the field
contributes to
\(\bra{\Omega}A_\mu(x)\ket{\mathbf k,\mathrm{null}}\).
The sum of longitudinal and time-like polarization vectors is proportional
to the null momentum:
\[
 \chi_\mu(x)=C\,k_\mu e^{ik\cdot x}.
\]
Choose
\[
 \lambda(x)=\frac{C}{i}e^{ik\cdot x}.
\]
Then \(\chi_\mu=\partial_\mu\lambda\), while
\[
 \Box\lambda=-k^2\lambda=0.
\]
It is precisely a residual gauge transformation preserving
\(\partial_\mu A^\mu=0\).  The null state therefore changes a potential
representative but produces no field strength and no physical photon.
\item
Let
\[
 \ket f=\int d^3k\,f(\mathbf k)
 \ket{\mathbf k,\mathrm{null}},
\]
where \(f\) is square integrable and sufficiently regular.  Linearity of
the field matrix element gives
\[
 \chi_\mu^{(f)}(x)
 =\int d^3k\,C(\mathbf k)f(\mathbf k)
 k_\mu e^{ik\cdot x}.
\]
Define
\[
 \lambda_f(x)=\int d^3k\,
 \frac{C(\mathbf k)f(\mathbf k)}{i}e^{ik\cdot x}.
\]
Every momentum in the integral is null, so
\(\Box\lambda_f=0\), and differentiating under the integral gives
\(\chi_\mu^{(f)}=\partial_\mu\lambda_f\).  Superposing residual gauge
transformations therefore never creates a physical transverse component.
\end{enumerate}

\SupplementarySolution{8}{Invariant Flux and Two-Body Phase Space}
\begin{enumerate}
\item
For mostly-plus momenta \(p_i^2=-m_i^2\), the Lorentz scalar
\[
 \mathcal F=4\sqrt{(p_1\cdot p_2)^2-m_1^2m_2^2}
\]
is the invariant flux.  In the center-of-mass frame,
\(p_1=(E_1,\mathbf p)\) and \(p_2=(E_2,-\mathbf p)\).
Using \(E_i^2=\mathbf p^2+m_i^2\) gives
\[
 \sqrt{(p_1\cdot p_2)^2-m_1^2m_2^2}
 =|\mathbf p|\,(E_1+E_2)=|\mathbf p|\sqrt s,
\]
so \(\mathcal F=4|\mathbf p|\sqrt s\).  In the rest frame of particle
2, \(p_2=(m_2,\mathbf0)\), and the square root becomes
\(m_2|\mathbf p_1|\).  Thus \(\mathcal F=4m_2|\mathbf p_1|\), the
familiar target-density times incident-velocity normalization in invariant
form.
\item
The invariant measure is
\[
 d\Phi_2=(2\pi)^4\delta^4(P-k_1-k_2)
 \frac{d^3k_1}{(2\pi)^3\,2E_1'}
 \frac{d^3k_2}{(2\pi)^3\,2E_2'}.
\]
In the center-of-mass frame the spatial delta function sets
\(\mathbf k_2=-\mathbf k_1\equiv-\mathbf k\).  The remaining delta
function is \(\delta(\sqrt s-E_1'(k)-E_2'(k))\).  Since
\[
 \frac{d(E_1'+E_2')}{dk}
 =k\left(\frac1{E_1'}+\frac1{E_2'}\right)
 =\frac{k\sqrt s}{E_1'E_2'},
\]
the radial integral gives
\[
 d\Phi_2=\frac{|\mathbf k|}{16\pi^2\sqrt s}\,d\Omega.
\]
Combining this with the flux
\(4|\mathbf p|\sqrt s\) yields, for distinguishable final particles,
\[
 \boxed{\;
 \frac{d\sigma}{d\Omega}
 =\frac{|\mathcal M|^2}{64\pi^2s}
 \frac{|\mathbf k|}{|\mathbf p|}.\;}
\]
\item
The invariant phase-space integral labels the variables
\(\mathbf k_1,\ldots,\mathbf k_n\).  If the particles are identical,
permuting these labels leaves the physical final state unchanged but
generates \(n!\) points in the labeled integration domain.  Division by
\(n!\) removes the overcounting.

For Møller scattering, the point with one electron at angle \(\theta\) and
the other at \(\pi-\theta\) is identical to the point obtained by exchanging
the two momenta.  One may therefore either integrate the labeled formula
over one hemisphere without an extra factor, or integrate over the full
solid angle and multiply by \(1/2!\).  Doing both would undercount, while
doing neither would double the physical rate.
\end{enumerate}

\SupplementarySolution{9}{Relativistic Rutherford Scattering}
\begin{enumerate}
\item
The first Born result is
\[
 f(\mathbf q)=-\frac{\mu_{\rm red}}{2\pi}\widetilde V(\mathbf q),
\qquad \frac{d\sigma}{d\Omega}=|f(\mathbf q)|^2.
\]
\item
For \(V(r)=Z\alpha/r\),
\(\widetilde V=4\pi Z\alpha/\mathbf q^2\) and
\(\mathbf q^2=4p^2\sin^2(\theta/2)\), yielding
\[
 \frac{d\sigma_R}{d\Omega}
 =\frac{Z^2\alpha^2}
 {4\mu_{\rm red}^2v^4\sin^4(\theta/2)}.
\]
\item
For a static source of charge \(Ze\),
\[
 A^0(\mathbf q)=\frac{Ze}{\mathbf q^2},\qquad
 \mathcal M=-e\,\bar u(p')\gamma^0u(p)\,A^0(\mathbf q).
\]
The spin average is evaluated with a trace:
\[
 \frac12\sum|\bar u(p')\gamma^0u(p)|^2
 =\frac12\operatorname{tr}
 [(\slashed p'+m)\gamma^0(\slashed p+m)\gamma^0].
\]
\item
Using \(|\mathbf p'|=|\mathbf p|\) and
\(\mathbf q^2=4|\mathbf p|^2\sin^2(\theta/2)\), the external-field
phase-space normalization gives the Mott formula
\[
 \boxed{\;
 \frac{d\sigma}{d\Omega}
 =\frac{Z^2\alpha^2}
 {4|\mathbf p|^2\beta^2\sin^4(\theta/2)}
 \left[1-\beta^2\sin^2\frac\theta2\right].\;}
\]
The bracket is the relativistic spin correction.  It tends to one for
\(\beta\ll1\), recovering Rutherford scattering.
\end{enumerate}

\SupplementarySolution{10}{Electron--Muon Scattering Invariants and Heavy-Target Limit}
\begin{enumerate}
\item
With \(p_1,p_2\) initial and \(k_1,k_2\) final electron and muon momenta,
the one \(t\)-channel graph gives
\[
 \mathcal M=\frac{e^2}{t}\,
 \bar u(k_1)\gamma^\mu u(p_1)\,
 \bar u(k_2)\gamma_\mu u(p_2),
\qquad t=-(p_1-k_1)^2.
\]
Two independent four-gamma traces yield
\[
 \boxed{\;
 \overline{|\mathcal M|^2}
 =\frac{2e^4}{t^2}\left[
 (s-m_e^2-m_\mu^2)^2
 +(u-m_e^2-m_\mu^2)^2
 +2t(m_e^2+m_\mu^2)
 \right].\;}
\]
Here
\(s+t+u=2m_e^2+2m_\mu^2\).  Setting both masses to zero gives the
standard check \(2e^4(s^2+u^2)/t^2\).
\item
In the center-of-mass frame,
\[
 p_*=\frac{\lambda^{1/2}(s,m_e^2,m_\mu^2)}{2\sqrt s},
\qquad
 \frac{d\sigma}{d\Omega}
 =\frac{\overline{|\mathcal M|^2}}{64\pi^2s},
\]
because elastic unequal-mass scattering has the same incoming and
outgoing \(p_*\).
\item
In the muon rest frame, let the incident electron have energy \(E\),
momentum \(p\), and speed \(\beta=p/E\).  For \(m_\mu\to\infty\),
the energy transfer to the target is \(O(1/m_\mu)\), so
\[
 E'=E,\qquad |\mathbf p'|=p,\qquad
 t=-4p^2\sin^2\frac\theta2.
\]
Substituting
\[
 s=m_\mu^2+m_e^2+2m_\mu E+O(1),\qquad
 u=m_\mu^2+m_e^2-2m_\mu E-t+O(1)
\]
into S.8.10(a) cancels the target-mass factors against the laboratory flux.
The remaining expression is
\[
 \frac{d\sigma}{d\Omega}
 =\frac{\alpha^2}
 {4p^2\beta^2\sin^4(\theta/2)}
 \left(1-\beta^2\sin^2\frac\theta2\right),
\]
the \(Z=1\) Mott--Rutherford result.  Infinite target mass removes recoil
but retains the electron's relativistic spin factor.
\end{enumerate}

\SupplementarySolution{11}{Electric Charge from a Static Potential and Dirac Magnetic Moment}
\begin{enumerate}
\item
For a weak classical field, the interaction matrix element is
\[
 \bra{p,s}H_I\ket{p',s'}
 =q\int d^3x\,e^{-i(\mathbf p-\mathbf p')\cdot\mathbf x}
 \bar u_s(p)\gamma^\mu u_{s'}(p')A_\mu(\mathbf x).
\]
Set \(\mathbf A=0\) and take \(|\mathbf p|,|\mathbf p'|\ll m\).
The temporal bilinear is
\[
 \bar u_s(p)\gamma^0u_{s'}(p')
 =2m\,\xi_s^\dagger\xi_{s'}+O(p^2/m).
\]
After dividing by the \(2m\) that converts relativistic to
nonrelativistic state normalization, the Hamiltonian is
\[
 H_{\rm NR}=qA^0(\mathbf x).
\]
Thus the charge measured by a static potential is exactly the coefficient
\(q\) in the covariant derivative.
\item
For a rigid sphere with identical normalized charge and mass profiles,
the rotational integrals differ only by \(q/m\), and the current carries
one additional factor \(1/2\):
\[
 \boldsymbol\mu=\frac{q}{2m}\mathbf L.
\]
Hence a classical rotating distribution has \(g=1\).
\item
The Gordon identity gives
\[
 \bar u(p')\gamma^\mu u(p)
 =\bar u(p')\left[
 \frac{(p'+p)^\mu}{2m}
 +\frac{i\sigma^{\mu\nu}(p'-p)_\nu}{2m}
 \right]u(p).
\]
For a static magnetic field, the first term supplies orbital minimal
coupling.  Fourier transforming the second and using
\(\mathbf B=\boldsymbol\nabla\times\mathbf A\) gives
\[
 H_{\rm spin}=-\frac{q}{m}\,\mathbf S\cdot\mathbf B.
\]
\item
Comparing with
\(-\boldsymbol\mu\cdot\mathbf B\), where
\(\boldsymbol\mu=gq\mathbf S/(2m)\), yields
\[
 \boxed{g=2}
\]
at tree level.  The difference from the classical rigid-body value
\(g=1\) follows from the spinor current rather than a literal rotating
charge distribution.
\end{enumerate}

\SupplementarySolution{12}{Nonminimal Electromagnetic Couplings}
For compact coefficients define the volume-convention Pauli tensor
\[
 \sigma^{\mu\nu}\equiv\frac{i}{2}[\gamma^\mu,\gamma^\nu].
\]
Thus the operator in part (a) is proportional to
\(\bar\psi\sigma^{\mu\nu}\psi F_{\mu\nu}\).  Overall phases needed to make
each Hamiltonian Hermitian are understood below; they are fixed by the
volume's anti-Hermitian vector gamma matrices and do not change the
nonrelativistic interpretation.
\begin{enumerate}
\item
Choose the conventional normalization
\[
 \Delta\mathcal L
 =-\frac{\kappa q}{4m}\,
 \bar\psi\sigma^{\mu\nu}\psi F_{\mu\nu}.
\]
For slowly moving positive-energy spinors, the spatial tensor bilinear
reduces to the Pauli spin density, while
\(F_{ij}\) is the magnetic field.  The resulting Hamiltonian term is
\[
 \Delta H=-\frac{\kappa q}{m}\,
 \mathbf S\cdot\mathbf B.
\]
Adding the minimal Dirac term gives
\[
 \boldsymbol\mu=\frac{q}{m}(1+\kappa)\mathbf S,
\qquad
 g=2(1+\kappa).
\]
Thus \(\kappa=(g-2)/2\) is the anomalous magnetic moment in this
normalization.  Gauge invariance permits the operator because it contains
the invariant field strength.
\item
With a real dipole coefficient \(d\), a Hermitian convention is
\[
 \Delta\mathcal L_{\rm EDM}
 =-\frac{i d}{2}\,
 \bar\psi\sigma^{\mu\nu}\gamma_5\psi F_{\mu\nu}.
\]
The dominant nonrelativistic bilinear converts the time--space components
of \(F_{\mu\nu}\) into
\[
 \Delta H_{\rm EDM}=-d\,\boldsymbol\sigma\cdot\mathbf E
 =-2d\,\mathbf S\cdot\mathbf E,
\]
up to whether \(d\) is defined for \(\boldsymbol\sigma\) or \(\mathbf S\).
Spin is an axial vector and \(\mathbf E\) a polar vector, so their dot
product is odd under parity.  Spin reverses under time reversal while
\(\mathbf E\) does not, making it time-reversal odd as well.  A permanent
EDM therefore signals \(P\) and \(T\), and under the usual assumptions
\(CP\), violation.
\item
Write the operator as
\[
 \Delta\mathcal L
 =\frac{c_V}{\Lambda^2}
 \bar\psi\gamma^\mu\psi\,\partial^\nu F_{\mu\nu}.
\]
For a classical external source, Maxwell's equation gives
\(\partial^\nu F_{\mu\nu}=j^{\rm ext}_\mu\), so
\[
 \Delta\mathcal L
 =\frac{c_V}{\Lambda^2}
 (\bar\psi\gamma^\mu\psi)j^{\rm ext}_\mu .
\]
It is a local current--current interaction, not a long-range dipole field.
In momentum space the photon vertex is proportional to
\(q^2\gamma^\mu-q^\mu\slashed q\).  Against a conserved source the second
term vanishes, leaving a contribution proportional to \(q^2\).  It
therefore changes the slope of the electric form factor at \(q^2=0\), often
interpreted as a charge-radius term, while vanishing for an on-shell real
photon.
\item
The axial analogue is
\[
 \Delta\mathcal L_A
 =\frac{a}{\Lambda^2}
 \bar\psi\gamma^\mu\gamma_5\psi\,
 \partial^\nu F_{\mu\nu}
 =\frac{a}{\Lambda^2}
 (\bar\psi\gamma^\mu\gamma_5\psi)j^{\rm ext}_\mu .
\]
For a slowly moving fermion, the spatial axial current is proportional to
the spin, so the leading term couples \(\mathbf S\) to the external
electric current rather than directly to \(\mathbf E\) or \(\mathbf B\).
The momentum-space vertex contains
\((q^2\gamma^\mu-q^\mu\slashed q)\gamma_5\), which vanishes for an
on-shell transverse photon.  It is therefore neither an electric nor a
magnetic dipole and is observable only through sources or virtual photons.
The anapole interaction is parity odd and time-reversal even; it is the
standard electromagnetic moment associated with parity violation without
a permanent EDM.
\end{enumerate}
\end{enumerate}
